\documentclass[11pt,a4paper]{article}
\usepackage[utf8]{inputenc}
\usepackage[margin=0.75in]{geometry}
\usepackage{amsmath, amssymb}
\usepackage{enumitem}
\usepackage{multicol}
\usepackage{titlesec}

\titleformat{\section}{\large\bfseries}{\thesection.}{0.5em}{}

\title{\textbf{NEET Biology: High-Difficulty Practice Set}\\ \large Topic: Mitosis and Cytokinesis}
\author{}
\date{}

\begin{document}

\maketitle

\section*{Questions}

\begin{enumerate}[label=\textbf{\arabic*.}, leftmargin=*]

\item A diploid cell ($2n = 24$) undergoes mitosis. If the DNA content in the $G_1$ phase is $40\text{ pg}$, what will be the number of chromosomes, chromatids, and DNA content in a single daughter nucleus immediately after telophase, assuming non-disjunction of one pair of homologous chromosomes occurred?
\begin{enumerate}[label=(\alph*)]
    \item $24$ chromosomes, $24$ chromatids, $40\text{ pg}$
    \item $25$ chromosomes, $25$ chromatids, $41.67\text{ pg}$
    \item $24$ chromosomes, $48$ chromatids, $80\text{ pg}$
    \item $22$ chromosomes, $22$ chromatids, $36.67\text{ pg}$
\end{enumerate}

\item Which of the following events marks the exact transition from early prophase to late prophase (prometaphase) in a typical eukaryotic animal cell?
\begin{enumerate}[label=(\alph*)]
    \item Initiation of assembly of mitotic spindle and condensation of chromatin
    \item Complete disintegration of the nuclear envelope and attachment of spindle fibers to kinetochores
    \item Movement of duplicated centrosomes to opposite poles and disappearance of nucleolus
    \item Splitting of centromeres and migration of daughter chromatids toward opposite poles
\end{enumerate}

\item \textbf{Assertion (A):} Colchicine induces polyploidy in dividing plant cells by arresting mitosis at metaphase.\\
\textbf{Reason (R):} Colchicine binds to tubulin dimers, preventing their polymerization into microtubules, thereby inhibiting sister chromatid separation while DNA replication remains unaffected.
\begin{enumerate}[label=(\alph*)]
    \item Both (A) and (R) are true and (R) is the correct explanation of (A).
    \item Both (A) and (R) are true but (R) is not the correct explanation of (A).
    \item (A) is true but (R) is false.
    \item Both (A) and (R) are false.
\end{enumerate}

\item During anaphase of mitosis, the centromere of each chromosome faces the pole while the arms trail behind. Which of the following chromosome morphologies will appear 'J-shaped' and 'L-shaped', respectively, during this stage?
\begin{enumerate}[label=(\alph*)]
    \item Metacentric and Sub-metacentric
    \item Acrocentric and Sub-metacentric
    \item Sub-metacentric and Acrocentric
    \item Telocentric and Metacentric
\end{enumerate}

\item Match Column I with Column II regarding mitotic inhibitors and regulators, and select the correct option:
\begin{center}
\begin{tabular}{|l|l|}
\hline
\textbf{Column I} & \textbf{Column II} \\
\hline
A. Taxol & 1. Inhibits cytokinesis by blocking actin polymerization \\
B. Cytochalasin B & 2. Stabilizes microtubules and prevents depolymerization \\
C. Nocodazole & 3. Induces chromosome condensation and nuclear envelope breakdown \\
D. MPF (Maturation Promoting Factor) & 4. Depolymerizes microtubules and arrests cells in prometaphase \\
\hline
\end{tabular}
\end{center}
\begin{enumerate}[label=(\alph*)]
    \item A-2, B-1, C-4, D-3
    \item A-4, B-1, C-2, D-3
    \item A-2, B-3, C-4, D-1
    \item A-1, B-2, C-3, D-4
\end{enumerate}

\item In an actively dividing meristematic tissue of an angiosperm, a cell is observed to have $32$ chromosomes at metaphase. How many kinetochores and spindle fibers (kinetochore microtubules) are associated with these chromosomes at this stage?
\begin{enumerate}[label=(\alph*)]
    \item $32$ kinetochores, $32$ spindle fibers
    \item $64$ kinetochores, $64$ spindle fibers
    \item $32$ kinetochores, $64$ spindle fibers
    \item $64$ kinetochores, $32$ spindle fibers
\end{enumerate}

\item Consider the following statements regarding cytokinesis in plant and animal cells:
\begin{enumerate}[label=(\roman*)]
    \item In animal cells, furrow formation occurs centripetally due to the contraction of a ring composed of actin and myosin II filaments.
    \item In plant cells, cell plate formation initiates at the periphery and grows inward (centripetally) to fuse with the existing lateral walls.
    \item Phragmoplast is a complex assembly of microtubules, microfilaments, and endoplasmic reticulum vesicles derived from the Golgi apparatus.
    \item Syncytium formation in liquid endosperm of coconut results from karyokinesis not followed by cytokinesis.
\end{enumerate}
Which of the above statements are correct?
\begin{enumerate}[label=(\alph*)]
    \item (i), (iii), and (iv) only
    \item (i), (ii), and (iii) only
    \item (ii) and (iv) only
    \item (i) and (iv) only
\end{enumerate}

\item What would be the immediate consequence if a human somatic cell in metaphase is treated with a drug that specifically inhibits the enzyme \textit{separase}?
\begin{enumerate}[label=(\alph*)]
    \item Chromosomes would fail to align at the equatorial plate.
    \item Cohesin complexes holding sister chromatids together would not be cleaved, preventing anaphase onset.
    \item Condensin proteins would fail to supercoil chromatin fibers during prophase.
    \item Nuclear envelope reformation would occur prematurely around unseparated chromatids.
\end{enumerate}

\item A root tip cell of \textit{Allium cepa} ($2n = 16$) undergoes $4$ successive mitotic divisions. How many total mitotic divisions (cycles) and how many total cells will be produced from a single parent cell, and what will be the chromosome number in each resulting cell?
\begin{enumerate}[label=(\alph*)]
    \item $15$ divisions, $16$ cells, $16$ chromosomes
    \item $4$ divisions, $16$ cells, $16$ chromosomes
    \item $16$ divisions, $16$ cells, $8$ chromosomes
    \item $15$ divisions, $32$ cells, $16$ chromosomes
\end{enumerate}

\item Identify the incorrect statement regarding the significance of mitosis:
\begin{enumerate}[label=(\alph*)]
    \item Mitosis restores the nucleo-cytoplasmic ratio that gets disturbed due to cell growth.
    \item In lower plants and social insects like male honeybees, haploid cells divide exclusively by mitosis.
    \item Mitosis introduces genetic recombination through sister chromatid exchange, aiding in evolution.
    \item Continuous mitotic divisions in the apical and lateral meristems result in indeterminate growth in plants.
\end{enumerate}

\item During which stage of mitosis does the chromosome reach its maximum level of condensation, making it ideal for karyotyping?
\begin{enumerate}[label=(\alph*)]
    \item Late Prophase
    \item Metaphase
    \item Early Anaphase
    \item Telophase
\end{enumerate}

\item Read the following events of mitosis:
\begin{enumerate}[label=(\Alph*)]
    \item Decondensation of chromosomes into chromatin network
    \item Splitting of centromeres
    \item Reformation of nucleolus, Golgi complex, and ER
    \item Alignment of chromosomes at the metaphase plate
    \item Movement of centrosomes to opposite poles
\end{enumerate}
Arrange them in the correct chronological sequence:
\begin{enumerate}[label=(\alph*)]
    \item E $\rightarrow$ D $\rightarrow$ B $\rightarrow$ A $\rightarrow$ C
    \item E $\rightarrow$ B $\rightarrow$ D $\rightarrow$ C $\rightarrow$ A
    \item D $\rightarrow$ E $\rightarrow$ B $\rightarrow$ A $\rightarrow$ C
    \item E $\rightarrow$ D $\rightarrow$ B $\rightarrow$ C $\rightarrow$ A
\end{enumerate}

\item In certain dinoflagellates and fungi, mitosis occurs within an intact nuclear envelope without the formation of a typical extranuclear spindle. This type of mitosis is termed:
\begin{enumerate}[label=(\alph*)]
    \item Eumitosis
    \item Cryptomitosis / Dinomitosis
    \item Endomitosis
    \item Amitosis
\end{enumerate}

\item If a cell has $2C$ DNA and $2n$ chromosomes in $G_1$, what is the ratio of the number of chromatids in metaphase to the number of chromosomes in anaphase within the same dividing cell?
\begin{enumerate}[label=(\alph*)]
    \item $1 : 1$
    \item $2 : 1$
    \item $1 : 2$
    \item $4 : 1$
\end{enumerate}

\item Which of the following structural features distinguishes plant mitosis from animal mitosis?
\begin{enumerate}[label=(\alph*)]
    \item Presence of kinetochores on centromeres
    \item Anastral spindle formation and centrifugal cytokinesis
    \item Amphiastral spindle formation and centripetal cytokinesis
    \item Disappearance of nuclear membrane at the end of prophase
\end{enumerate}

\item \textbf{Statement I:} During anaphase, Anaphase-Promoting Complex/Cyclosome (APC/C) degrades securin, which in turn activates separase.\\
\textbf{Statement II:} Anaphase A involves the shortening of kinetochore microtubules, whereas Anaphase B involves the sliding of polar microtubules past each other to elongate the spindle.
\begin{enumerate}[label=(\alph*)]
    \item Both Statement I and Statement II are correct.
    \item Both Statement I and Statement II are incorrect.
    \item Statement I is correct but Statement II is incorrect.
    \item Statement I is incorrect but Statement II is correct.
\end{enumerate}

\item A cell is treated with \textit{Cytochalasin B} during early prophase. After completion of the cell cycle, the resulting cell will be:
\begin{enumerate}[label=(\alph*)]
    \item Haploid and uninucleate
    \item Diploid and binucleate
    \item Tetraploid and uninucleate
    \item Arrested at metaphase
\end{enumerate}

\item At the end of telophase, which of the following organelles do NOT reform because they are semi-autonomous and divide by fission prior to cytokinesis?
\begin{enumerate}[label=(\alph*)]
    \item Golgi apparatus and Endoplasmic Reticulum
    \item Nucleolus and Nuclear envelope
    \item Mitochondria and Chloroplasts
    \item Both (a) and (b)
\end{enumerate}

\item How many of the following structures are visible under a light microscope in a human somatic cell at the end of prophase?
\textit{[Golgi bodies, Endoplasmic reticulum, Nucleolus, Nuclear envelope, Condensed chromosomes, Centrosomes]}
\begin{enumerate}[label=(\alph*)]
    \item Two
    \item Three
    \item Four
    \item Five
\end{enumerate}

\item During metaphase, chromosomes are held at the equator by balanced bipolar tension. What happens if one kinetochore of a chromosome fails to attach to spindle fibers from its respective pole?
\begin{enumerate}[label=(\alph*)]
    \item The cell immediately enters anaphase, resulting in one aneuploid daughter cell.
    \item The Spindle Assembly Checkpoint (SAC) generates a 'wait' signal, inhibiting APC/C activation.
    \item The unattached chromosome undergoes programmed degradation by lysosomes.
    \item Both sister chromatids are pulled to the same pole without centromere splitting.
\end{enumerate}

\item Match the sub-stages/events of M-phase with their molecular mechanisms:
\begin{center}
\begin{tabular}{|l|l|}
\hline
\textbf{Event} & \textbf{Molecular Mechanism} \\
\hline
A. Nuclear lamina breakdown & 1. Dephosphorylation of nuclear lamins \\
B. Chromosome condensation & 2. Phosphorylation of nuclear lamins by Cdk1-Cyclin B \\
C. Nuclear envelope reassembly & 3. Activation of Condensin complexes \\
D. Sister chromatid separation & 4. Cleavage of Scc1 subunit of cohesin \\
\hline
\end{tabular}
\end{center}
\begin{enumerate}[label=(\alph*)]
    \item A-2, B-3, C-1, D-4
    \item A-1, B-3, C-2, D-4
    \item A-2, B-4, C-1, D-3
    \item A-3, B-2, C-1, D-4
\end{enumerate}

\item In an organism with $2n = 14$, what is the total number of telomeres present in a cell at mitotic metaphase?
\begin{enumerate}[label=(\alph*)]
    \item $14$
    \item $28$
    \item $56$
    \item $112$
\end{enumerate}

\item Which of the following best describes the chemical nature and origin of the cell plate in higher plants?
\begin{enumerate}[label=(\alph*)]
    \item Cellulose and hemicellulose secreted by mitochondria
    \item Calcium and magnesium pectate formed by fusion of phragmoplast vesicles from Golgi
    \item Suberin and lignin deposited by smooth endoplasmic reticulum
    \item Chitin and peptidoglycan synthesized by plasma membrane enzymes
\end{enumerate}

\item \textbf{Assertion (A):} In animal cells, cytokinesis is achieved by the appearance of a furrow in the plasma membrane.\\
\textbf{Reason (R):} The position of the cleavage furrow is determined by the mitotic spindle, specifically the overlapping astral and interpolar microtubules at the equator.
\begin{enumerate}[label=(\alph*)]
    \item Both (A) and (R) are true and (R) is the correct explanation of (A).
    \item Both (A) and (R) are true but (R) is not the correct explanation of (A).
    \item (A) is true but (R) is false.
    \item Both (A) and (R) are false.
\end{enumerate}

\item What is the key difference between \textit{Endomitosis} and \textit{Free-nuclear division}?
\begin{enumerate}[label=(\alph*)]
    \item Endomitosis occurs without spindle formation and nuclear membrane breakdown, leading to polyploidy within a single nucleus; free-nuclear division involves karyokinesis without cytokinesis, producing multinucleate cells.
    \item Endomitosis produces haploid cells; free-nuclear division produces diploid cells.
    \item Endomitosis occurs only in animals; free-nuclear division occurs only in bacteria.
    \item Endomitosis involves cytokinesis without karyokinesis; free-nuclear division involves both simultaneously.
\end{enumerate}

\item A diploid cell ($2n = 20$) has $10\text{ pg}$ of DNA in its $G_1$ phase. If this cell undergoes mitosis, what will be the amount of DNA per chromosome during metaphase and anaphase, respectively?
\begin{enumerate}[label=(\alph*)]
    \item $1.0\text{ pg}$ and $0.5\text{ pg}$
    \item $0.5\text{ pg}$ and $0.5\text{ pg}$
    \item $1.0\text{ pg}$ and $1.0\text{ pg}$
    \item $2.0\text{ pg}$ and $1.0\text{ pg}$
\end{enumerate}

\item During telophase, chromosomes cluster at opposite spindle poles and lose their individuality. Which of the following processes does \textbf{NOT} occur during telophase?
\begin{enumerate}[label=(\alph*)]
    \item Disassembly of the mitotic spindle apparatus
    \item Reassembly of the nuclear envelope around each chromosome cluster
    \item Replication of centrioles for the next cell cycle
    \item Decondensation of chromosomes into euchromatin and heterochromatin
\end{enumerate}

\item Which of the following types of microtubules in the mitotic spindle do NOT attach to chromosomes but overlap at the spindle midzone to push the poles apart?
\begin{enumerate}[label=(\alph*)]
    \item Kinetochore microtubules
    \item Astral microtubules
    \item Polar (interpolar) microtubules
    \item Cortical microtubules
\end{enumerate}

\item In a hypothetical experiment, a dividing animal cell is microinjected with antibodies against \textit{myosin II} during late anaphase. What will be the observed phenotype?
\begin{enumerate}[label=(\alph*)]
    \item Sister chromatids will fail to move to opposite poles.
    \item Nuclear envelope will not reform around daughter nuclei.
    \item Contractile ring will assemble properly, allowing normal cleavage.
    \item Cleavage furrow formation will fail, resulting in a binucleate cell.
\end{enumerate}

\item Identify the correct statement regarding the kinetochore:
\begin{enumerate}[label=(\alph*)]
    \item It is a non-coding DNA sequence located at the ends of chromosomes.
    \item It is a disc-shaped proteinaceous structure on the surface of the centromere that serves as the attachment site for spindle microtubules.
    \item Each chromosome possesses only one kinetochore during metaphase.
    \item It is composed exclusively of actin and myosin proteins.
\end{enumerate}

\item Why is mitosis referred to as 'equational division' even though the DNA content per cell is halved during M-phase compared to the $G_2$ phase?
\begin{enumerate}[label=(\alph*)]
    \item Because the chromosome number and DNA content of the daughter cells are identical to those of the parent cell in the $G_1$ phase.
    \item Because crossing over occurs equally between all homologous chromosomes.
    \item Because the volume of cytoplasm is divided in an exact $1:1$ ratio in all cells.
    \item Because both karyokinesis and cytokinesis take equal time to complete.
\end{enumerate}

\item Observe the following features of a dividing cell:
\begin{itemize}
    \item Spindle fibers are attached to kinetochores of chromosomes.
    \item Chromosomes are moved to spindle equator and get aligned along the metaphase plate.
    \item Each chromosome consists of two sister chromatids held together at the centromere.
\end{itemize}
If the ploidy of this cell is $4n = 48$, how many centromeres and chromatids are present at this stage?
\begin{enumerate}[label=(\alph*)]
    \item $48$ centromeres, $96$ chromatids
    \item $24$ centromeres, $48$ chromatids
    \item $48$ centromeres, $48$ chromatids
    \item $96$ centromeres, $96$ chromatids
\end{enumerate}

\item In higher plants, the preprophase band (PPB) is a transient ring of microtubules and actin filaments that forms just beneath the plasma membrane before mitosis. What is its primary function?
\begin{enumerate}[label=(\alph*)]
    \item It forms the mitotic spindle in the absence of centrosomes.
    \item It marks the exact cortical division site where the new cell plate will fuse with the parental wall during cytokinesis.
    \item It induces condensation of chromatin fibers during early prophase.
    \item It degrades the middle lamella to allow cell expansion.
\end{enumerate}

\item Which of the following motor proteins is primarily responsible for the poleward movement of chromosomes along kinetochore microtubules during Anaphase A?
\begin{enumerate}[label=(\alph*)]
    \item Cytoplasmic dynein and depolymerizing kinesins (e.g., Kinesin-13)
    \item Myosin V
    \item Axonemal dynein
    \item Kinesin-1 moving toward the minus end
\end{enumerate}

\item \textbf{Statement I:} In animal cells, centrosome duplication occurs during the S phase, while their migration to opposite poles occurs during prophase.\\
\textbf{Statement II:} Radiating microtubules around each centrosome during prophase are called asters, and two asters together with spindle fibers form the mitotic apparatus.
\begin{enumerate}[label=(\alph*)]
    \item Both Statement I and Statement II are correct.
    \item Both Statement I and Statement II are incorrect.
    \item Statement I is correct but Statement II is incorrect.
    \item Statement I is incorrect but Statement II is correct.
\end{enumerate}

\item A mutant yeast strain has a temperature-sensitive defect in the gene encoding \textit{Cyclin B}. When shifted to the restrictive temperature during $G_2$, the cells will:
\begin{enumerate}[label=(\alph*)]
    \item Arrest at the $G_2$/M transition and fail to enter prophase.
    \item Enter mitosis normally but arrest at the metaphase-anaphase transition.
    \item Undergo uncontrolled multiple rounds of mitosis without S phase.
    \item Complete karyokinesis but fail to undergo cytokinesis.
\end{enumerate}

\item Which of the following events does \textbf{NOT} happen during prophase of mitosis?
\begin{enumerate}[label=(\alph*)]
    \item Untangling of chromosomal material during chromatin condensation
    \item Movement of duplicated centrioles to opposite poles of the cell
    \item Attachment of spindle fibers to kinetochores of chromosomes
    \item Disappearance of Golgi complexes, ER, nucleolus, and nuclear envelope at the end of the phase
\end{enumerate}

\item During anaphase, as sister chromatids separate and move toward opposite poles, what happens to the distance between the centromeres of sister chromatids and the distance between the centromeres and their respective poles?
\begin{enumerate}[label=(\alph*)]
    \item Both distances increase.
    \item Both distances decrease.
    \item Distance between sister centromeres increases; distance between centromere and pole decreases.
    \item Distance between sister centromeres decreases; distance between centromere and pole increases.
\end{enumerate}

\item How many mitotic generations are required to produce $128$ cells from a single meristematic cell?
\begin{enumerate}[label=(\alph*)]
    \item $7$
    \item $64$
    \item $127$
    \item $8$
\end{enumerate}

\item In certain protozoans like \textit{Amoeba} and in bacterial cells, direct cell division occurs without spindle formation or distinct chromosome condensation, where the nucleus simply constricts and divides into two. This is called:
\begin{enumerate}[label=(\alph*)]
    \item Mitosis
    \item Meiosis
    \item Amitosis
    \item Karyokinesis
\end{enumerate}

\item If the $G_1$ phase of a diploid cell has $20$ chromosomes and $2C$ DNA content, which of the following rows correctly represents the state at Anaphase?
\begin{center}
\begin{tabular}{|l|c|c|c|}
\hline
\textbf{Option} & \textbf{Chromosomes per cell} & \textbf{Chromatids per cell} & \textbf{DNA content per cell} \\
\hline
(a) & $20$ & $40$ & $4C$ \\
(b) & $40$ & $40$ & $4C$ \\
(c) & $40$ & $0$ & $2C$ \\
(d) & $20$ & $20$ & $2C$ \\
\hline
\end{tabular}
\end{center}

\item Which of the following is the correct sequence of layers formed during plant cytokinesis from the center outward?
\begin{enumerate}[label=(\alph*)]
    \item Primary cell wall $\rightarrow$ Middle lamella $\rightarrow$ Plasma membrane
    \item Middle lamella $\rightarrow$ Primary cell wall $\rightarrow$ Plasma membrane
    \item Plasma membrane $\rightarrow$ Primary cell wall $\rightarrow$ Middle lamella
    \item Secondary cell wall $\rightarrow$ Primary cell wall $\rightarrow$ Middle lamella
\end{enumerate}

\item \textbf{Assertion (A):} Syncytium is a multinucleate condition that arises when karyokinesis is not followed by cytokinesis.\\
\textbf{Reason (R):} The cellular endosperm in coconut is a classic example of syncytium.
\begin{enumerate}[label=(\alph*)]
    \item Both (A) and (R) are true and (R) is the correct explanation of (A).
    \item Both (A) and (R) are true but (R) is not the correct explanation of (A).
    \item (A) is true but (R) is false.
    \item Both (A) and (R) are false.
\end{enumerate}

\item During mitotic prometaphase, chromosomes exhibit congression. What does 'congression' refer to?
\begin{enumerate}[label=(\alph*)]
    \item Condensation of chromatin into visible chromosomes
    \item Movement of chromosomes toward the equatorial plane of the spindle
    \item Splitting of centromeres at the onset of anaphase
    \item Clustering of chromosomes at opposite poles during telophase
\end{enumerate}

\item Consider a cell in mitotic telophase. If a scientist measures the amount of histone proteins and DNA in one of the two newly formed daughter nuclei and compares it to the parent cell in $G_2$ phase, the ratio (Daughter nucleus at telophase : Parent cell at $G_2$) will be:
\begin{enumerate}[label=(\alph*)]
    \item $1 : 1$
    \item $1 : 2$
    \item $2 : 1$
    \item $1 : 4$
\end{enumerate}

\end{enumerate}

\newpage
\section*{Answer Key}

\begin{multicols}{5}
\begin{enumerate}[label=\textbf{\arabic*.}]
    \item (b)
    \item (b)
    \item (a)
    \item (b)
    \item (a)
    \item (b)
    \item (a)
    \item (b)
    \item (a)
    \item (c)
    \item (b)
    \item (a)
    \item (b)
    \item (a)
    \item (b)
    \item (a)
    \item (b)
    \item (c)
    \item (a)
    \item (b)
    \item (a)
    \item (c)
    \item (b)
    \item (a)
    \item (a)
    \item (a)
    \item (c)
    \item (c)
    \item (d)
    \item (b)
    \item (a)
    \item (a)
    \item (b)
    \item (a)
    \item (a)
    \item (a)
    \item (c)
    \item (c)
    \item (a)
    \item (c)
    \item (b)
    \item (b)
    \item (c)
    \item (b)
    \item (b)
\end{enumerate}
\end{multicols}

\section*{Detailed Solutions}

\begin{enumerate}[label=\textbf{\arabic*.}, leftmargin=*]

\item \textbf{(b)} Normally, a diploid cell ($2n = 24$, DNA $= 40\text{ pg}$ in $G_1$) duplicates DNA in S phase to $80\text{ pg}$ and divides equally into two daughter cells with $24$ chromosomes and $40\text{ pg}$ DNA each. If non-disjunction of one chromosome occurs (note: in mitosis, sister chromatids of one chromosome fail to separate, resulting in $2n+1$ and $2n-1$), one daughter cell receives $24 + 1 = 25$ chromosomes. Since DNA per chromatid is $40/24 = 1.67\text{ pg}$, $25$ chromatids equal $40 + 1.67 = 41.67\text{ pg}$.

\item \textbf{(b)} According to NCERT, the complete disintegration of the nuclear envelope marks the start of the second phase of mitosis (metaphase/prometaphase), allowing chromosomes to spread through the cytoplasm and attach to spindle fibers via kinetochores.

\item \textbf{(a)} Colchicine acts as a mitotic poison by binding to tubulin subunits and preventing microtubule assembly. Without a functional spindle, sister chromatids cannot separate at anaphase, leading to chromosome doubling (polyploidy) within the same cell.

\item \textbf{(b)} During anaphase, centromeres lead while arms trail. Metacentric chromosomes appear \textbf{V-shaped}, sub-metacentric appear \textbf{L-shaped}, acrocentric appear \textbf{J-shaped}, and telocentric appear \textbf{I-shaped}.

\item \textbf{(a)} Taxol hyper-stabilizes microtubules (used in chemotherapy); Cytochalasin B blocks actin polymerization and inhibits cleavage furrow formation; Nocodazole depolymerizes microtubules; MPF (Cdk1 + Cyclin B) drives entry into M-phase by triggering chromatin condensation and nuclear envelope breakdown.

\item \textbf{(b)} At metaphase, each of the $32$ chromosomes consists of $2$ sister chromatids, and each chromatid has its own kinetochore on the centromere. Therefore, there are $32 \times 2 = 64$ kinetochores, and each attaches to at least one spindle fiber ($64$ spindle fibers).

\item \textbf{(a)} Statement (ii) is incorrect because in plant cells, cell plate formation initiates at the center of the cell and grows outward (\textbf{centrifugally}) to meet the existing lateral walls. All other statements are correct.

\item \textbf{(b)} During the metaphase-to-anaphase transition, the enzyme \textit{separase} cleaves the cohesin rings holding sister chromatids together. Inhibiting separase prevents sister chromatid separation and blocks anaphase onset.

\item \textbf{(a)} Number of cells produced after $n$ mitotic generations $= 2^n = 2^4 = 16$ cells. However, the total number of individual mitotic divisions required to produce $N$ cells from a single parent cell is $N - 1 = 16 - 1 = 15$ divisions. Mitosis is equational, so chromosome number remains $2n = 16$.

\item \textbf{(c)} Mitosis produces genetically identical daughter cells and does \textbf{not} involve crossing over or genetic recombination (which occurs in meiosis).

\item \textbf{(b)} Condensation of chromosomes is completed by metaphase, making them shortest and thickest, which is ideal for studying chromosome morphology and karyotyping.

\item \textbf{(a)} Correct chronological order: Centrosomes move to opposite poles (Prophase, E) $\rightarrow$ Chromosomes align at equator (Metaphase, D) $\rightarrow$ Splitting of centromeres (Anaphase, B) $\rightarrow$ Decondensation of chromosomes (Telophase, A) $\rightarrow$ Reformation of nucleolus, Golgi, and ER (Late Telophase, C).

\item \textbf{(b)} Intranuclear mitosis where the nuclear membrane remains intact throughout division is called \textit{cryptomitosis} or \textit{dinomitosis} (seen in dinoflagellates and many fungi).

\item \textbf{(a)} In a $2n$ cell, metaphase has $2n$ chromosomes, each with $2$ chromatids $= 4n$ chromatids. In anaphase, centromeres split, so each chromatid becomes an independent chromosome $= 4n$ chromosomes. Thus, the ratio is $4n : 4n = 1 : 1$.

\item \textbf{(b)} Higher plant cells lack centrioles, forming an \textit{anastral} spindle, and divide their cytoplasm via a cell plate that grows \textit{centrifugally} (center to periphery). Animal cells have an \textit{amphiastral} spindle and divide \textit{centripetally} (furrowing).

\item \textbf{(a)} Both statements are accurate molecular descriptions of anaphase: APC/C ubiquitinates securin, freeing separase to cleave cohesin; Anaphase A pulls chromatids via kinetochore microtubule shortening, while Anaphase B pushes poles apart via polar microtubules.

\item \textbf{(b)} Cytochalasin B inhibits actin microfilaments, blocking cytokinesis without affecting karyokinesis. Thus, nuclear division completes normally, yielding a single diploid cell with two diploid nuclei (binucleate).

\item \textbf{(c)} Mitochondria and chloroplasts are semi-autonomous organelles with their own DNA; they divide by fission during interphase and are distributed between daughter cells rather than disassembling and reforming like the nuclear envelope, ER, and Golgi.

\item \textbf{(a)} By the end of prophase, the Golgi complex, ER, nucleolus, and nuclear envelope disappear. Only the condensed chromosomes and centrosomes (asters) remain visible.

\item \textbf{(b)} The Spindle Assembly Checkpoint (SAC) monitors bipolar attachment of all kinetochores. Even a single unattached kinetochore generates a signal (Mad2/MCC) that inhibits APC/C, halting progression to anaphase until proper attachment occurs.

\item \textbf{(a)} Cdk1-Cyclin B phosphorylates nuclear lamins to disassemble the nuclear lamina (A-2); condensin supercoils DNA (B-3); phosphatases dephosphorylate lamins to rebuild the nuclear envelope (C-1); separase cleaves the Scc1 subunit of cohesin (D-4).

\item \textbf{(c)} In a cell with $2n = 14$ at metaphase, there are $14$ duplicated chromosomes, meaning $28$ chromatids. Each linear chromatid has $2$ ends (telomeres). Therefore, total telomeres $= 28 \times 2 = 56$.

\item \textbf{(b)} The cell plate is formed by the coalescence of Golgi-derived vesicles guided by the phragmoplast, which deposit calcium and magnesium pectate to form the middle lamella.

\item \textbf{(a)} Both statements are true, and (R) correctly explains (A): signals from the overlapping interpolar and astral microtubules at the spindle midzone (via RhoA pathway) specify where the actin-myosin contractile ring forms.

\item \textbf{(a)} Endomitosis is chromosome replication and chromatid separation inside an intact nucleus without spindle or cytokinesis, creating polyploid nuclei (e.g., megakaryocytes, tapetum). Free-nuclear division involves repeated karyokinesis without cytokinesis, creating a coenocyte/syncytium.

\item \textbf{(a)} In $G_1$, $20$ chromosomes have $10\text{ pg}$ DNA ($0.5\text{ pg}$ per chromatid). In metaphase, each chromosome has $2$ chromatids $= 1.0\text{ pg}$ DNA per chromosome. In anaphase, sister chromatids separate into individual chromosomes, so each chromosome has $1$ chromatid $= 0.5\text{ pg}$ DNA.

\item \textbf{(c)} Centriole replication occurs during the **S phase** of interphase, not during telophase.

\item \textbf{(c)} Polar (interpolar) microtubules extend from opposite poles and overlap at the equatorial midzone, where kinesin-5 motor proteins slide them past each other to elongate the spindle during Anaphase B.

\item \textbf{(d)} Myosin II provides the motor force for contracting the actin ring during animal cytokinesis. Blocking Myosin II prevents furrow ingression, leading to a binucleate cell.

\item \textbf{(b)} Kinetochores are trilaminar disc-shaped protein structures assembled on the centromeric region of each sister chromatid (so $2$ per metaphase chromosome) where spindle microtubules attach.

\item \textbf{(a)} Although DNA is halved from $4C$ (in $G_2$) to $2C$ (in daughter cells), each daughter cell ends up with the exact same chromosome number ($2n$) and DNA content ($2C$) as the parent cell had before entering S phase ($G_1$).

\item \textbf{(a)} Regardless of ploidy ($4n = 48$), at metaphase there are $48$ chromosomes, meaning $48$ centromeres. Since each chromosome is duplicated into two sister chromatids, there are $48 \times 2 = 96$ chromatids.

\item \textbf{(b)} The preprophase band (PPB) predicts and establishes the cortical division site in plant cells; even though it disassembles before metaphase, it leaves behind molecular markers that guide the expanding phragmoplast during telophase.

\item \textbf{(a)} Poleward movement during Anaphase A is driven by minus-end-directed motor proteins (cytoplasmic dynein) at the kinetochores coupled with microtubule depolymerization by Kinesin-13 (MCAK).

\item \textbf{(a)} Both statements are direct facts from NCERT Class 11 Biology (Chapter: Cell Cycle and Cell Division).

\item \textbf{(a)} Cyclin B is required to activate Cdk1 (forming MPF) to trigger entry into mitosis. Without functional Cyclin B, cells cannot pass the $G_2$/M checkpoint.

\item \textbf{(c)} Attachment of spindle fibers to kinetochores occurs during **metaphase** (specifically prometaphase/early metaphase) after the complete disintegration of the nuclear envelope, whereas events (a), (b), and (d) occur during prophase.

\item \textbf{(c)} As sister chromatids are pulled toward opposite poles in anaphase, they move away from each other (increasing the distance between sister centromeres) and closer to their respective spindle poles (decreasing centromere-to-pole distance).

\item \textbf{(a)} Each mitotic generation doubles the cell number: $2^n = 128 \implies 2^n = 2^7 \implies n = 7$ generations.

\item \textbf{(c)} Amitosis is direct nuclear and cytoplasmic cleavage without chromosome condensation or spindle formation, seen in prokaryotes and certain protozoans/aging tissues.

\item \textbf{(b)} At Anaphase, centromeres split, doubling the chromosome count per cell from $20$ to $40$. Since each separated chromatid is now counted as an independent chromosome (or single-chromatid unit), there are $40$ chromosomes (and $40$ chromatid units) sharing the $4C$ DNA within the single dividing cell prior to cytokinesis.

\item \textbf{(b)} The cell plate matures into the central \textbf{middle lamella}, on either side of which each daughter cell deposits its \textbf{primary cell wall}, bounded internally by the \textbf{plasma membrane}.

\item \textbf{(c)} Assertion (A) is true, but Reason (R) is false because it is the **liquid endosperm** (free-nuclear endosperm) of coconut that is a syncytium, not the cellular endosperm (white kernel).

\item \textbf{(b)} Congression is the active movement and alignment of chromosomes onto the metaphase plate (equatorial plane) via opposing pulling forces from both spindle poles.

\item \textbf{(b)} A parent cell in $G_2$ has $4C$ DNA and twice the normal histone content ($2\times$). After mitosis, each daughter nucleus at telophase receives exactly half of the DNA ($2C$) and half of the histones ($1\times$), giving a ratio of $1 : 2$.

\end{enumerate}

\end{document}